% Copyright (c) 2026 Angela Villota, and collaborators from the CyED block
% Licensed under the PolyForm Noncommercial License 1.0.0.
% Commercial use is prohibited without prior written authorization.
%===============================================================================
% LaTeX  Diapositivas Computaci�n y Estructuras Discretas III 2023-2 ICESI
% Andr��s A. Aristiz�bal P. - �ngela P. Villota G.
%===============================================================================
%===============================================================================
\documentclass[compress,xcolor={dvipsnames}]{beamer}
\mode<presentation> {
%============================================================================== 
\usetheme{Boadilla}
%\usetheme{Copenhagen}
%\usetheme{CambridgeUS}
%============================================================================== 
%\usecolortheme{crane}
%\usecolortheme{seahorse}
%============================================================================== 
%\usecolortheme{crane}
\usecolortheme{whale}
%============================================================================== 
\usefonttheme[onlylarge]{structurebold}

%==============================================================================
%==============================================================================
\setbeamertemplate{navigation symbols}{}
\setbeamertemplate{footline}[frame number]
\setbeamerfont*{frametitle}{size=\normalsize,series=\bfseries}
%\setbeamercolor{structure}{fg=black!80!white,bg=white}
%\setbeamercolor{palette quaternary}{fg=white!80!gray,bg=red!50!black}}

\setbeamercolor{structure}{fg=white,fg=blue!60!black}  %%%%Estaba en 70
\setbeamercolor{palette quaternary}{fg=gray!40!black,bg=gray!20!white}}


%=============================================================================== 
%========================================================================================
\usepackage[spanish,activeacute]{babel}
\usepackage[applemac]{inputenc}
\usepackage[spanish]{babel}
\usepackage{times}
\usepackage{latexsym}
\usepackage{color}
\usepackage{amsmath}
\usepackage{amsfonts}
\usepackage{amssymb}
\usepackage{stmaryrd}
\usepackage{indentfirst}
\usepackage{multicol}
\usepackage{graphicx}
\usepackage{graphics}
\usepackage{newlfont}
\usepackage{exscale}
\usepackage[T1]{fontenc}

\usepackage{tikz}
\usepackage{amscd}
\usepackage{pb-diagram}
%\usepackage{pstricks}
\usepackage{makeidx}
\usepackage{textpos}

%==============================================================================
%==============================================================================
\tikzstyle{block}=[draw opacity=0.7,line width=1.4cm]
%============================================================================== 
%=======================================================================================
\newtheorem{teo}{Teorema}
\newtheorem{prop}{Proposition}
\newtheorem{afir}{Afirmaci�n}
\newtheorem{obs}{Observaci�n}
\newtheorem{lem}{Lemma}
\newtheorem{cor}{Corollary}
\newtheorem{ej}{Example}
\newtheorem{ejer}{Exercise}
\newtheorem{ejers}{Exercises}
\newtheorem{ejs}{Examples}
\newtheorem{df}{Definition}

%%%%%%%%%%%%%%%%%%%%%%%%%%%%%%%%%%%%%%%%%%%%%%%%%%%%%%%%%%%%%%%%%%%%%Flechas
\newcommand{\stateequiv}{\,\sim\,}
\newcommand{\entonces}{\,\rightarrow\,}
\newcommand{\pertenece}{\,\in \,}
\newcommand{\talque}{\,| \,}
\newcommand{\sii}{\,\leftrightarrow\,}
\newcommand{\ra}{\,\rightarrow\,}
\newcommand{\lra}{\leftrightarrow}
\newcommand{\Ra}{\;\Longrightarrow\;}
\newcommand{\lora}{\longrightarrow}
\newcommand{\we}{\,\wedge\,}
\newcommand{\ve}{\,\vee \,}
\newcommand{\bbsq}{\blacksquare\!\blacksquare}
%%%%%%%%%%%%%%%%%%%%%%%%%%%%%%%%%%%%%%%%%%%%%%%%%%%%%%%%%%%%%%%%%%%%%S�mbolos l�gicos
\newcommand{\mo}{\models}
\newcommand{\lan}{\langle}
\newcommand{\ran}{\rangle}
\newcommand{\ap}{\approx}
%%%%%%%%%%%%%%%%%%%%%%%%%%%%%%%%%%%%%%%%%%%%%%%%%%%%%%%%%%%%%%%%%%%%%Letras
\newcommand{\vep}{\varepsilon}
\newcommand{\vph}{\varphi}
\newcommand{\de}{\delta}
\newcommand{\si}{\sigma}
\newcommand{\Si}{\Sigma}
\newcommand{\om}{\omega}
\newcommand{\al}{\alpha}
\newcommand{\be}{\beta}
\newcommand{\vom}{\varOmega}
\newcommand{\vTh}{\varTheta}
\newcommand{\vth}{\vartheta}
\newcommand{\vDe}{\varDelta}
\newcommand{\wh}{\widehat}
\newcommand{\nn}{\~n}
\newcommand{\m}{i}
\newcommand{\ps}{P_{\Sigma}}
\newcommand{\psone}{P_{\Sigma^{*}}}
\newcommand{\psoneast}{P_{\Sigma'}}
\newcommand{\ts}{T_{\Sigma}}
\newcommand{\dis}{\displaystyle}
\newenvironment{variableblock}[3]{%
  \setbeamercolor{block body}{#2}
  \setbeamercolor{block title}{#3}
  \begin{block}{#1}}{\end{block}}

%============================================================================== 
%==========================================================================================

\title{Computaci\'on y Estructuras Discretas III}

\author[Andr�s A. Aristiz�bal P.]{Andr�s A. Aristiz�bal P. aaaristizabal@icesi.edu.co \\ Angela Villota - apvillota@icesi.edu.co}

\institute[Universidad ICESI]{Departamento de Computaci\'on y Sistemas Inteligentes \\
\includegraphics[height=1.5cm]{../logoICESI2.pdf}}

\date[]{\today}

%==============================================================================
\AtBeginSubsection[]{

%%%%%%%%%%%%%%%%%%%%%%%%%%%%%%%%%%%%%%%%%%%%%%%%%%%%%%%%%%%%%%%%%%%%%%%%%%%%%%
\begin{frame}
    \frametitle{Agenda}
    \tableofcontents[currentsubsection, 
    hideothersubsections, 
    sectionstyle=show/shaded, 
    subsectionstyle=show/shaded,]
  \end{frame}}
%%%%%%%%%%%%%%%%%%%%%%%%%%%%%%%%%%%%%%%%%%%%%%%%%%%%%%%%%%%%%%%%%%%%%%%%%%%%%%

%============================================================================== CUERPO DEL DOCUMENTO
%%%%%%%%%%%%%%%%%%%%%%%%%%%%%%%%%%%%%%%%%%%%%%%%%%%%%%%%%%%%%%%%%%%%%%%%%%%%%%
\begin{document}
\maketitle
\setlength{\parindent}{0cm}
%%%%%%%%%%%%%%%%%%%%%%%%%%%%%%%%%%%%%%%%%%%%%%%%%%%%%%%%%%%%%%%%%%%%%%%%%%%%%%
%\begin{frame}
 %   \titlepage
%\end{frame}
%%%%%%%%%%%%%%%%%%%%%%%%%%%%%%%%%%%%%%%%%%%%%%%%%%%%%%%%%%%%%%%%%%%%%%%%%%%%%%

%%%%%%%%%%%%%%%%%%%%%%%%%%%%%%%%%%%%%%%%%%%%%%%%%%%%%%%%%%%%%%%%%%%%%%%%%%%%%%
\begin{frame}
\frametitle{Agenda del d�a}
  \tableofcontents
\end{frame}
%%%%%%%%%%%%%%%%%%%%%%%%%%%%%%%%%%%%%%%%%%%%%%%%%%%%%%%%%%%%%%%%%%%%%%%%%%%%%%

%%%%%%%%%%%%%%%%%%%%%%%%%%%%%%%%%%%%%%%%%%%%%%%%%%%%%%%%%%%%%%%%%%%%%%%%%%%%%%

\setbeamertemplate{itemize/enumerate body begin}{\footnotesize}

%\section[Regular languages and Automata theory]{Regular languages and Automata theory}
%\section{Operations Using Languages}
%
%%%%%%%%%%%%%%
%{\footnotesize\begin{frame}{Operations as Sets}
%
%\begin{df}
%Since languages over $\Sigma$ are subsets of $\Sigma^*,$ the usual operations between sets are also valid operations between languages. Then, if $A$ and $B$ are languages over $\Sigma$ ($A, \, B \subseteq \Sigma^*$), then the following are also languages over $\Sigma:$
%\begin{center}
%\begin{tabular}{ll}
%$A \cup B$ & Union \\
%$A \cap B$ & Intersection \\
%$A-B$ & Difference \\
%$\overline{A} = \Sigma^*-A$ & Complement 
%\end{tabular}
%\end{center}
%\end{df}
%These operations between languages are called boolean operations in order to differentiate them from the linguistic operations, which are extensions to the languages of the operations between strings.
%\end{frame}}
%%%%%%%%%%%%%%
%\subsection{Concatenation}
%\begin{frame}{Concatenation}
%
%\begin{df}
%The concatenation of two languages $A$ and $B$ over $\Sigma,$ represented as $A \cdot B$ or simply $A B$ is defined as: \pause
%\[A B = \{u v \talque u \pertenece A, \, v \pertenece B\}\] \pause
%In general, $A B \neq B A$
%\end{df}
%\end{frame}
%
%\setbeamertemplate{itemize/enumerate body begin}{\footnotesize}
%
%%%%%%%%%%%%%%
%{\scriptsize \begin{frame}{Properties over Concatenation}
%
%\begin{df}
%Given languages $A, B, C$ over $\Sigma,$ i.e.  $A, \, B, \, C \subseteq \Sigma^*.$ Then \pause
%\begin{enumerate}
%\item $A \cdot \emptyset = \emptyset \cdot A = \emptyset.$ 
%\item $A \cdot \{\lambda\} = \{\lambda\} \cdot A = A.$ 
%\item Associative property,
%\begin{center}
%$A \cdot (B \cdot C) = (A \cdot B) \cdot C.$
%\end{center} 
%\item Distributivity of concatenation with respect to union,
%\begin{center}
%$A \cdot (B \cup C) = A \cdot B \cup A \cdot C.$ \\
%$(B \cup C) \cdot A = B \cdot A \cup C \cdot A.$ 
%\end{center} 
%\item Generalized distributivity property. If $\{B_i\}_{i \pertenece I}$ is any family of languages over $\Sigma,$ then
%\begin{center}
%$A \cdot \bigcup\limits_{i \pertenece I} B_i = \bigcup\limits_{i \pertenece I} (A \cdot B_i),$ \\
%$\Big( \bigcup\limits_{i \pertenece I} B_i\Big) \cdot A = \bigcup\limits_{i \pertenece I} (B_i \cdot A).$ 
%\end{center} 
%\end{enumerate}
%\end{df}
%\end{frame}}
%
%
%%\subsection{Exercises}
%%%%%%%%%%%%%%%%%%%%%%%%%%%%%%%%%
%%{\begin{frame}{Exercises}
%%\begin{ejer}
%%If $u,v \pertenece \Sigma^*,$ show that $(u v)^R = v^R u^R.$ 
%%\end{ejer}
%%\end{frame}}
%%%%%%%%%%%%%%%%%%%%%%%%%%%%%%%%%
%%{\begin{frame}{Exercises}
%%\begin{ejer}
%%Show that $A \cdot (B \cdot C) = (A \cdot B) \cdot C.$
%%\end{ejer}
%%\end{frame}}
%%%%%%%%%%%%%%%%%%%%%%%%%%%%%%%%%
%%{\begin{frame}{Exercises}
%%\begin{ejer}
%%Show that $A \cdot \{\lambda\} = \{\lambda\} \cdot A = A.$
%%\end{ejer}
%%\end{frame}}
%%%%%%%%%%%%%%%%%%%%%%%%%%%%%%%%%%
%%{\begin{frame}{Exercises}
%%\begin{ejer}
%%Show that $A \cdot \emptyset = \emptyset \cdot A = \emptyset.$
%%\end{ejer}
%%\end{frame}}
%%%%%%%%%%%%%%%%%%%%%%%%%%%%%%%%%%%%%%%%%%%%%%%%%%%%%%%%%%%%%%%%%%%%%%%%%%%%%%%

\subsection{Power}
{\footnotesize \begin{frame}{Power}
What is the power of a language? 
\begin{df}
Given a language $A$ over $\Sigma,$ $(A \subseteq \Sigma^*),$ and a natural number $n \pertenece \mathbb{N},$ $A^n$ is defined in the following way: 
\begin{center}
\begin{tabular}{lll}
$A^0$ & $=$ & $\{\lambda\},$ \\ 
$A^n$ & $=$ & $\underbrace{A A \cdot \cdot \cdot A}_\text{n times} = \{u_1 \cdot \cdot \cdot u_n \talque u_i \pertenece A, \text{ for all } i, \, 1 \leq i \leq n \}.$ 
\end{tabular}
\end{center}
\end{df}
$A^2$ is the set of double string concatenations of $A,$ $A^3$ the set of triple string concatenations of $A$ and in general $A^n$ is the set of $n$ string concatenations of $A.$
\end{frame}}
\subsection{Kleene Closure}
{\footnotesize\begin{frame}{Kleene Closure}
\begin{df}
The Kleene closure of $A,$ $A \subseteq \Sigma^*,$ is the union of all powers of $A$ and it is represented by $A^*.$ 
\begin{center}
\begin{tabular}{lll}
$A^*$ & $=$ & set of all string concatenations of $A,$ \\ 
& & including $\lambda$ \\ 
 & $=$ & $\{u_1, \cdot \cdot \cdot u_n \talque u_i \pertenece A, \, n \geq 0\}$ 
\end{tabular}
\end{center}
\end{df}

\end{frame}}

{\footnotesize\begin{frame}{Positive Closure}

\begin{df}
The positive closure of a language $A,$ $A \subseteq \Sigma^*,$ is the union of all powers of $A,$ without including $\lambda$ and is represented by $A^+.$ 
\begin{center}
\begin{tabular}{lll}
$A^+$ & $=$ & set of all string concatenations of $A,$ \\ 
& & without including $\lambda$ \\ 
 & $=$ & $\{u_1, \cdot \cdot \cdot u_n \talque u_i \pertenece A, \, n \geq 1\}$ 
\end{tabular}
\end{center}
\end{df}

\end{frame}}
%%%%%%%%%%%%
{\footnotesize\begin{frame}{Properties}
Which are the properties of $*$ and $+$? 
\begin{df}
\begin{enumerate}
\item $A^+ = A^* \cdot A = A \cdot A^*.$ 
\item $A^* \cdot A^* = A^*.$ 
\item ($A^*)^n = A^*,$ for all n, $n \geq 1.$  
\item ($A^*)^* = A^*.$ 
\item $A^+ \cdot A^+ \subseteq A^+.$  
\item ($A^*)^+ = A^*.$  
\item ($A^+)^* = A^*.$  
\item ($A^+)^+ = A^+.$ 
\item If $A$ and $B$ are languages over $\Sigma^*,$ then $(A \cup B)^* = (A^* B^*)^*.$
\end{enumerate}
\end{df}
\end{frame}}
\subsection{Reverse}
{\footnotesize\begin{frame}{Reverse}

\begin{df}
Given $A,$ a language over $\Sigma,$ we define $A^R$ in the following way:
\[A^R = \{u^R \talque u \pertenece A\}\] 
\end{df}
$A^R$ is the inverse of $A.$
\end{frame}}

%
%%%%%%%%%%%%%
%%{\footnotesize\begin{frame}{Languages}
%%Which are the properties of the inverse of a language? \pause
%%\begin{df}
%%Given $A$ and $B$ langauges over $\Sigma$ ($A, \, B \subseteq \Sigma^*).$ \pause
%%\begin{enumerate}
%%\item $(A \cdot B)^R = B^R \cdot A^R.$ \pause
%%\item $(A \cup B)^R = A^R \cup B^R.$ \pause
%%\item $(A \cap B)^R = A^R \cap B^R.$ \pause
%%\item $(A^R)^R = A.$ \pause
%%\item $(A^*)^R = (A^R)^*.$ \pause
%%\item $(A^+)^R = (A^R)^+.$ 
%%\end{enumerate}
%%\end{df}
%%\end{frame}}
%%%%%%%%%%%%%
%

\subsection{Languages}
{\footnotesize \begin{frame}{Languages}
What is the power of a language? \pause
\begin{df}
Given a language $A$ over $\Sigma,$ $(A \subseteq \Sigma^*),$ and a natural numberl $n \pertenece \mathbb{N},$ $A^n$ is defined in the following way: \pause
\begin{center}
\begin{tabular}{lll}
$A^0$ & $=$ & $\{\lambda\},$ \\ 
$A^n$ & $=$ & $\underbrace{A A \cdot \cdot \cdot A}_\text{n times} = \{u_1 \cdot \cdot \cdot u_n \talque u_i \pertenece A, \text{ for all } i, \, 1 \leq i \leq n \}.$ \pause
\end{tabular}
\end{center}
\end{df}
$A^2$ is the set of double string concatenations of $A,$ $A^3$ the set of triple string concatenations of $A$ and in general $A^n$ is the set of $n$ string concatenations of $A.$
\end{frame}}

{\footnotesize\begin{frame}{Languages}
What is the Kleene closure of a language? \pause
\begin{df}
The Kleene closure of $A,$ $A \subseteq \Sigma^*,$ is the union of all powers of $A$ and it is represented by $A^*.$ \pause
\begin{center}
\begin{tabular}{lll}
$A^*$ & $=$ & set of all string concatenations of $A,$ \\ 
& & including $\lambda$ \\ 
 & $=$ & $\{u_1, \cdot \cdot \cdot u_n \talque u_i \pertenece A, \, n \geq 0\}$ 
\end{tabular}
\end{center}
\end{df}

\end{frame}}

{\footnotesize\begin{frame}{Languages}
What is the positive closure of a language? \pause
\begin{df}
The positive closure of a language $A,$ $A \subseteq \Sigma^*,$ is the union of all powers of $A,$ without including $\lambda$ and is represented by $A^+.$ \pause
\begin{center}
\begin{tabular}{lll}
$A^+$ & $=$ & set of all string concatenations of $A,$ \\ 
& & without including $\lambda$ \\ 
 & $=$ & $\{u_1, \cdot \cdot \cdot u_n \talque u_i \pertenece A, \, n \geq 1\}$ 
\end{tabular}
\end{center}
\end{df}

\end{frame}}
%%%%%%%%%%%%
{\footnotesize\begin{frame}{Languages}
Which are the properties of $*$ and $+$? \pause
\begin{df}
\begin{enumerate}
\item $A^+ = A^* \cdot A = A \cdot A^*.$ \pause
\item $A^* \cdot A^* = A^*.$ \pause
\item ($A^*)^n = A^*,$ para todo $n \geq 1.$  \pause
\item ($A^*)^* = A^*.$  \pause
\item $A^+ \cdot A^+ \subseteq A^+.$  \pause
\item ($A^*)^+ = A^*.$  \pause
\item ($A^+)^* = A^*.$  \pause
\item ($A^+)^+ = A^+.$ \pause
\item If $A$ and $B$ are languages over $\Sigma^*,$ then $(A \cup B)^* = (A^* B^*)^*.$
\end{enumerate}
\end{df}
\end{frame}}

{\footnotesize\begin{frame}{Languages}
What is the inverse of a language? \pause
\begin{df}
Given $A,$ a language over $\Sigma,$ we define $A^R$ in the following way:
\[A^R = \{u^R \talque u \pertenece A\}\] \pause
\end{df}
$A^R$ is the inverse of $A.$
\end{frame}}


%%%%%%%%%%%%
{\footnotesize\begin{frame}{Languages}
Which are the properties of the inverse of a language? \pause
\begin{df}
Given $A$ and $B$ langauges over $\Sigma$ ($A, \, B \subseteq \Sigma^*).$ \pause
\begin{enumerate}
\item $(A \cdot B)^R = B^R \cdot A^R.$ \pause
\item $(A \cup B)^R = A^R \cup B^R.$ \pause
\item $(A \cap B)^R = A^R \cap B^R.$ \pause
\item $(A^R)^R = A.$ \pause
\item $(A^*)^R = (A^R)^*.$ \pause
\item $(A^+)^R = (A^R)^+.$ 
\end{enumerate}
\end{df}
\end{frame}}
%%%%%%%%%%%%
{\footnotesize\begin{frame}{Languages}
What are regular languages? \pause
\begin{block}{}
Regular languages from a given alphabet $\Sigma$ are all those languages that can be built from the basic languages $\emptyset, \{\lambda\}, \{a\}, a \pertenece \Sigma,$ using union, concatenation and Kleene closure. \pause
\end{block}
Given an alphabet $\Sigma$ \pause
\begin{df}
\begin{enumerate}
\item $\emptyset, \{\lambda\}, \{a\},$ for all $a \pertenece \Sigma,$ are regular languages over $\Sigma.$ These are the so called basic regular languages. \pause
\item If $A$ and $B$ are regular languages over $\Sigma,$ then the following are also regular: \pause
\begin{center}
\begin{tabular}{cl}
$A \cup B$ & (union), \\ \pause
$A \cdot B$ & (concatenation), \\ \pause
$A^*$ & (Kleene closure). 
\end{tabular}
\end{center}
\end{enumerate}
\end{df}
\end{frame}}
%%%%%%%%%%%%
{\scriptsize \begin{frame}{Languages}
\begin{block}{}
$\Sigma$ and $\Sigma^*$ are both regular languages over $\Sigma.$ Union, concatenation and Kleene closure are known as regular operations.
\end{block} \pause
\begin{ej}
Given $\Sigma = \{a,b\}.$ The following ones are regular languages over $\Sigma:$ \pause
\begin{enumerate}
\item The language $A$ of all strings that have just one $a:$ 
\begin{center}
$A = \{b\}^* \cdot \{a\} \cdot \{b\}^*.$ \pause
\end{center}
\item The language $B$ of all strings that start with $b:$ 
\begin{center}
$B = \{b\} \cdot \{a,b\}^*.$ \pause
\end{center}
\item The language $C$ of all strings that contain the string $ba:$ 
\begin{center}
$C = \{a,b\}^* \cdot \{ba\} \cdot \{a,b\}^*.$ \pause
\end{center}
\item $(\{a\} \cup \{b\}^*) \cdot \{a\}.$ \pause 
\item $[(\{a\}^* \cup \{b\}^*) \cdot \{b\}]^*.$ 
\end{enumerate}
\end{ej}
\end{frame}}

\section{Regular languages and expressions}
\subsection{Regular Languages}
{\footnotesize\begin{frame}{Regular Languages}

\begin{block}{}
Regular languages from a given alphabet $\Sigma$ are all those languages that can be built from the basic languages $\emptyset, \{\lambda\}, \{a\}, a \pertenece \Sigma,$ using union, concatenation and Kleene closure. \pause
\end{block}
Given an alphabet $\Sigma$ 
\begin{df}
\begin{enumerate}
\item $\emptyset, \{\lambda\}, \{a\},$ for all $a \pertenece \Sigma,$ are regular languages over $\Sigma.$ These are the so called basic regular languages. 
\item If $A$ and $B$ are regular languages over $\Sigma,$ then the following are also regular: 
\begin{center}
\begin{tabular}{cl}
$A \cup B$ & (union), \\ 
$A \cdot B$ & (concatenation), \\ 
$A^*$ & (Kleene closure). 
\end{tabular}
\end{center}
\end{enumerate}
\end{df}
\end{frame}}
%%%%%%%%%%%%
{\scriptsize \begin{frame}{Regular Languages}
\begin{block}{}
$\Sigma$ and $\Sigma^*$ are both regular languages over $\Sigma.$ Union, concatenation and Kleene closure are known as regular operations.
\end{block} 
\begin{ej}
Given $\Sigma = \{a,b\}.$ The following ones are regular languages over $\Sigma:$ 
\begin{enumerate}
\item The language $A$ of all strings that have just one $a:$ 
\begin{center}
$A = \{b\}^* \cdot \{a\} \cdot \{b\}^*.$ 
\end{center}
\item The language $B$ of all strings that start with $b:$ 
\begin{center}
$B = \{b\} \cdot \{a,b\}^*.$ 
\end{center}
\item The language $C$ of all strings that contain the string $ba:$ 
\begin{center}
$C = \{a,b\}^* \cdot \{ba\} \cdot \{a,b\}^*.$ 
\end{center}
\item $(\{a\} \cup \{b\}^*) \cdot \{a\}.$ 
\item $[(\{a\}^* \cup \{b\}^*) \cdot \{b\}]^*.$ 
\end{enumerate}
\end{ej}
\end{frame}}
\subsection{Regular Expressions}
%%%%%%%%%%%%
{\footnotesize\begin{frame}{Regular Expressions}
\textbf{Regular expression} is a sequence of symbols and characters expressing a \textbf{Regular Language}, they are used as strings or patterns to be searched for within a longer piece of text.

\end{frame}}

{\footnotesize\begin{frame}{Regular Expressions}

\begin{df}
\begin{enumerate}
\item Basic regular expressions:
\begin{center}
\begin{tabular}{ll}
$\emptyset$ & is a regular expression that represents the language $\emptyset$  \\ 
$\lambda$ & is a regular expression that represents the language $\{\lambda\}.$  \\ 
$a$ & is a regular expression that represents the language $\{a\},$ \\
& $a \pertenece \Sigma.$ 
\end{tabular}
\end{center}
\item If $R$ and $S$ are regular expressions over $\Sigma,$ the following ones are also regular: 
\begin{center}
$(R)(S)$ \\ 
$(R \cup S)$ \\ 
$(R)^*$ 
\end{center}

\end{enumerate}
\end{df}

\end{frame}}
%%%%%%%%%%%%%
{\footnotesize\begin{frame}{Example}
\begin{ej}
Given the alphabet $\Sigma = \{a,b,c\},$ \\ 
\[(a \cup b^*) a^* (bc)^* \] 
is the regular expression that represents 
\[(\{a\} \cup \{b\}^*) \cdot \{a\}^* \cdot \{bc\}^* \] 
\end{ej}
\end{frame}}





\begin{frame}{An overview of regular expressions}
\textbf{In programming, a regular expression is }
\begin{block}{}
\begin{itemize}
\item A pattern of text that consists of \textbf{literals} (ordinary characters) and special characters known as \textbf{metacharacters}. 
\item e.g. \texttt{a\textbackslash w*}
\end{itemize}
\end{block}

\begin{block}{}
We can also find
\begin{itemize}
\item Character classes (\texttt{[..]})
\item Negation of ranges ($[ \textasciicircum  \dots]$)
\item Sets of regular expressions  ($|$)
\item Quantifiers
\item Boundary matchers
\end{itemize}
\end{block}

\end{frame}

\begin{frame}{An overview of regular expressions}
What are literals? 
\begin{block}{}
\begin{itemize}
\item Are the simplest form of pattern matching in regular expression.
\item They will simply succeed whenever that literal is found.
\item If we apply the regular expression \textbf{fox }to search the phrase T\textit{he quick brown fox jumps over the lazy dog}, we will find one match.
\item We can also obtain several results instead of just one, if we apply the
regular expression \texttt{be} to the following phrase To be, or not to be.
\end{itemize}
\end{block}
\end{frame}

\begin{frame}{An overview of regular expressions}
What are metacharacters? 
\begin{block}{}
\begin{itemize}
\item A metacharacter is a character that has a special meaning during pattern processing. 
\item We use metacharacters in regular expressions to define the search criteria and any text manipulations. 
\item The simplest example of a metacharacter is the full stop.
\item The full stop character matches any single character of any sort (apart from a newline).
\item e.g. \texttt{.at} means: any letter, followed by the letter \texttt{a} followed by the letter \texttt{t}.
\end{itemize}
\end{block}
\end{frame}

\begin{frame}{An overview of regular expressions}
What are character classes? 
\begin{block}{}
\begin{itemize}
\item The character classes (also known as character sets) allow us to define a character that will match if any of the defined characters on the set is present.
\item To define a character class, we should use the opening square bracket metacharacter [, then any accepted characters, and finally close with a closing square bracket ].
\item e.g. \texttt{licen[cs]e}
\item It is possible to also use the range of a character. 
\ This is done by leveraging the hyphen symbol (-) between two related characters.
\item e.g. To match any lowercase letter we can use \texttt{[a-z].} 
\item e.g. To match any single digit we can define the character set \texttt{[0-9].}
\end{itemize}
\end{block}
\end{frame}

\begin{frame}{An overview of regular expressions}
How to work with several character classes? 
\begin{block}{}
\begin{itemize}
\item The character classes' \, ranges can be combined to be able to match a character against many ranges. 
\item We just need to put one range after the other. 
\item e.g. If we want to match any lowercase or uppercase alphanumeric character, we can use \texttt{[0-9a-zA-Z]} 
\item This can be alternatively written using the union mechanism: \texttt{[0-9[a-z[A-Z]]].}
\end{itemize}
\end{block}
\end{frame}

\begin{frame}{An overview of regular expressions}
What do we mean by negation of ranges? 
\begin{block}{}
\begin{itemize}
\item We can invert the meaning of a character set by placing a caret (\textasciicircum) symbol right after the opening square bracket metacharacter ([). 
\item If we have a character class such as \texttt{[0-9]} meaning any digit, the negated character class \texttt{[\textasciicircum0-9]} will match anything that is not a digit. 
\end{itemize}
\end{block}
\end{frame}

%\begin{frame}{An overview of regular expressions}
%Some character classes supported at this moment in Python 
%\begin{block}{}
%\begin{itemize}
%\item \texttt{.} This element matches any character except newline \texttt{\textbackslash n} \pause
%\item \texttt{\textbackslash d} This matches any decimal digit. \pause
%\item \texttt{\textbackslash D} This matches any non-digit character. \pause
%\item \texttt{\textbackslash s} This matches any whitespace character. \pause
%\item \texttt{\textbackslash S} This matches any non whitespace character. \pause
%\item \texttt{\textbackslash w} This matches any alphanumeric character. \pause
%\item \texttt{\textbackslash W} This matches any non alphanumeric character. 
%\end{itemize}
%\end{block}
%\end{frame}

\begin{frame}{An overview of regular expressions}
How to match against a set of regular expressions? \pause
\begin{block}{}
\begin{itemize}
\item This is accomplished by using the pipe symbol \texttt{|}. \pause
\item It alternates between regular expressions. \pause
\item e.g. \texttt{yes|no} \pause
\item e.g. \texttt{yes|no|maybe}
\end{itemize}
\end{block}
\end{frame}

\begin{frame}{An overview of regular expressions}
What are quantifiers? \pause
\begin{block}{}
\begin{itemize}
\item The mechanisms to define how a character, metacharacter, or character set can be repeated. \pause
\end{itemize}
\end{block}
Which are the three basic quantifiers? \pause
\begin{block}{}
\begin{itemize}
\item \texttt{?} 0 or 1 repetitions \pause
\item \texttt{*} 0 or more times \pause
\item \texttt{+} 1 or more times \pause
\item \texttt{\{n,m\}} Between $n$ and $m$ times. \pause
\begin{itemize}
\item \texttt{\{n\}} Exactly $n$ times. \pause
\item \texttt{\{n,\}} $n$ or more times. \pause
\item \texttt{\{,n\}} At most $n$ times. \pause
\end{itemize}
\item e.g. \texttt{cars?} Singular or plural.
\end{itemize}
\end{block}
\end{frame}

\begin{frame}{An overview of regular expressions}
\begin{ej}
How to match a telephone number that can be in the format 555-555-555, 555 555 555, or 555555555. 
\end{ej}
\end{frame}

%\begin{frame}{An overview of regular expressions}
%What are greedy and reluctant quantifiers? \pause
%\begin{block}{}
%\begin{itemize}
%\item The greedy behavior of the quantifiers is applied by default in the quantifiers. \pause
%\item A greedy quantifier will try to match as much as possible to have the biggest match result possible. \pause
%\item The non-greedy behavior can be requested by adding an extra question mark to the quantifier. \pause
%\item For example, \texttt{??}, \texttt{*?} or \texttt{+?}. \pause
%\item A quantifier marked as reluctant will behave like the exact opposite of the greedy ones. \pause
%\item They will try to have the smallest match possible.
%\end{itemize}
%\end{block}
%\end{frame}

%\begin{frame}{An overview of regular expressions}
%\begin{ej}
%Given the string \texttt{English $\rightarrow$ "Hello", $\rightarrow$ Spanish $\rightarrow$ "Hola".} What would be the results of using the regular expressions \texttt{".+"} and \texttt{".+?"} over that string.
%\end{ej}
%\end{frame}

\begin{frame}{An overview of regular expressions}
What are boundary matchers? \pause
\begin{block}{}
\begin{itemize}
\item The boundary matchers are a set or characters that will correspond to a particular position inside of the input. \pause
\item \texttt{\textasciicircum} Matches at the beginning of a line \pause
\item \texttt{\$} Matches at the end of a line \pause
\item \texttt{\textbackslash b} Matches a word boundary \pause
\item \texttt{\textbackslash B} The opposite of  \textbackslash b \pause
\item \texttt{\textbackslash A} Matches the beginning of the input \pause
\item \texttt{\textbackslash Z} Matches the end of the input 
\end{itemize}
\end{block}
\end{frame}

\begin{frame}{An overview of regular expressions}
\begin{ej}
Write a regular expression that will match lines that start with \texttt{Name: } and make sure that after the name, there are only alphabetic characters or spaces until the end of the line.
\end{ej}
\end{frame}

\begin{frame}{An overview of regular expressions}
\begin{ej}
What is the difference between the regular expressions \texttt{hello} and \texttt{\textbackslash bhello\textbackslash b}?
\end{ej}
\end{frame}

\section{Regular expressions in Python}
\begin{frame}{Regular expressions in Python}
How are regular expressions supported in Python? 
\begin{block}{}
\begin{itemize}
\item They are supported by the \texttt{re} module.
\item We only need to import it to start using it. 
\end{itemize}
\end{block}
How to start using them to match a pattern? 
\begin{block}{}
\begin{itemize}
\item We compile a pattern with \texttt{pattern = re.compile(r'foo')} 
\item We can try to match it against a string \texttt{pattern.match("foo bar")}
\end{itemize}
\end{block}
\end{frame}

\begin{frame}{Regular expressions in Python}
What are the building blocks for Python Regex? 
\begin{block}{}
\begin{itemize}
\item RegexObject
\begin{itemize}
\item Also known as Pattern Object. 
\item Represents a compiled regular expression.
\end{itemize}
\item MatchObject
\begin{itemize}
\item Represents the matched pattern.
\end{itemize}
\end{itemize}
\end{block}
\end{frame}

\begin{frame}{Regular expressions in Python}
What is a RegexObject? 
\begin{block}{}
\begin{itemize}
\item Before matching patterns we need to compile the regex. 
\item The compilation produces a reusable pattern object. 
\item This object provides all the operations that can be done (i.e. matching a pattern and finding all substrings that match a particular regex). 
\item e.g. \texttt{pattern = re.compile(r`<HTML>')} 
\item \texttt{pattern.match(``<HTML>")}
\end{itemize}
\end{block}
\end{frame}

\begin{frame}{Regular expressions in Python}
What are two ways of matching a pattern? 
\begin{block}{}
\begin{itemize}
\item We can compile a pattern, which gives us a RegexObject. 
\item We can just use the module operations.
\item If we compile a pattern we are able to reuse it.
\item e.g. \texttt{pattern = re.compile(r`<HTML>')} 
\item \texttt{pattern.match(``<HTML>")}
\item e.g. \texttt{re.match(r`<HTML>', ``<HTML>")} 
\end{itemize}
\end{block}
\end{frame}

\begin{frame}{Regular expressions in Python}
How to search for string that match a pattern? 
\begin{block}{}
\begin{itemize}
\item In python we have two operations match and search. 
\end{itemize}
\end{block}
How does match work? 
\begin{block}{}
\begin{itemize}
\item This method tries to match the compiled pattern only at the beginning of the string. 
\item If there is a match, then it returns a MatchObject.
\item It has two optional parameters, pos and endpos.
\item pos determines the position from where to search the pattern in the string.
\item endpos determines the position until where the pattern is searched in the string.
\item \texttt{pattern.match(string,pos,endpos)}
\end{itemize}
\end{block}
\end{frame}

\begin{frame}{Regular expressions in Python}
\begin{ej}
Given \texttt{pattern = re.compile(r`\textasciicircum<HTML>')} what is the result of: \pause 
\begin{itemize}
\item \texttt{pattern.match(``<HTML>")} 
\item \texttt{pattern.match(``  <HTML>")} 
\item \texttt{pattern.match(``  <HTML>"[2:])}
\end{itemize}
\end{ej}
\end{frame}

\begin{frame}{Regular expressions in Python}
\begin{ej}
Given \texttt{pattern = re.compile(r`<HTML>\$')} what is the result of: 
\begin{itemize}
\item \texttt{pattern.match(``<HTML> ",0,6)} 
\item \texttt{pattern.match(``<HTML> "[:6])}
\end{itemize}
\end{ej}
\end{frame}

\begin{frame}{Regular expressions in Python}
How does search work? 
\begin{block}{}
\begin{itemize}
\item This operation would be like the match of many languages.
\item It tries to match the pattern at any location of the string and not just at the beginning. 
\item If there is a match, then it returns a MatchObject.
\item The pos and endpos parameters have the same meaning as that in the match
operation.
\item pos determines the position from where to search the pattern in the string.
\item endpos determines the position until where the pattern is searched in the string.
\item \texttt{pattern.match(string,pos,endpos)}
\end{itemize}
\end{block}
\end{frame}

\begin{frame}{Regular expressions in Python}
\begin{ej}
Given \texttt{pattern = re.compile(r``world'')} what is the result of: 
\begin{itemize}
\item \texttt{pattern.match(``hello world")} 
\item \texttt{pattern.match(``hola mundo '')}
\end{itemize}
\end{ej}
\end{frame}

\begin{frame}{Regular expressions in Python}
\begin{ej}
Given \texttt{pattern = re.compile(r`\textasciicircum<HTML>', re.MULTILINE)} what is the result of: 
\begin{itemize}
\item \texttt{pattern.search(``<HTML>")} 
\item \texttt{pattern.search(`` <HTML>'')} 
\item \texttt{pattern.search(`` \textbackslash n<HTML>'')} 
\item \texttt{pattern.search(``  \textbackslash n<HTML>'',3)} 
\item \texttt{pattern.search(``</div></body>\textbackslash n<HTML>'',4)} 
\item \texttt{pattern.search(``  \textbackslash n<HTML>'',4)}
\end{itemize}
\end{ej}
\end{frame}

\begin{frame}{Regular expressions in Python}
How does findall work? 
\begin{block}{}
\begin{itemize}
\item It returns a list with all the non-overlapping occurrences of a pattern and not the MatchObject like search and match do. 
\item e.g. \texttt{pattern = re.compile(r``\textbackslash w+")}
\item \texttt{pattern.findall(``hello world'')}
\item \texttt{[`hello', `world']}
\end{itemize}
\end{block}
\end{frame}

\begin{frame}{Regular expressions in Python}
\begin{ej}
\begin{itemize}
\item \texttt{pattern = re.compile(r"(\textbackslash w+) (\textbackslash w+)")}
\item \texttt{pattern.findall(``Hello world hola mundo'')}
\item \texttt{[(`Hello', `world'), (`hola', 'mundo')]}
\end{itemize}
\end{ej}
\end{frame}

\begin{frame}{Regular expressions in Python}
How does finditer work? 
\begin{block}{}
\begin{itemize}
\item Works essentially as as findall.
\item It returns an iterator in which each element is a MatchObject.
\item We can use the operations provided by this object. 
\item Useful when you need information for every match.
\end{itemize}
\end{block}
\end{frame}

\begin{frame}{Regular expressions in Python}
\begin{ej}
\begin{itemize}
\item \texttt{pattern = re.compile(r``(\textbackslash w+) (\textbackslash w+)")}
\item \texttt{it = pattern.finditer(``Hello world hola mundo'')}
\item \texttt{match = it.next()}
\item \texttt{match.groups()}
\item \texttt{(`Hello',`world')}
\item \texttt{match.span()}
\item \texttt{(0,11)}
\end{itemize}
\end{ej}
\end{frame}

\begin{frame}{Regular expressions in Python}
Which are some operations to modify strings? \pause
\begin{block}{}
\begin{itemize}
\item split(string, maxsplit=0): A string can be split based on the matches of the pattern.
\item sub(repl, string, count=0): Returns the resulting string after replacing the matched pattern in the original string with the replacement.
\end{itemize}
\end{block}
\end{frame}

\begin{frame}{Regular expressions in Python}
\begin{ej}
\begin{itemize}
\item \texttt{pattern = re.compile(r``\textbackslash W")}
\item \texttt{pattern.split(``Beautiful is better than ugly", 2)}
\item \texttt{[`Beautiful', `is', `better than ugly']}
\end{itemize}
\end{ej}
\end{frame}


\begin{frame}{Regular expressions in Python}
\begin{ej}
\begin{itemize}
\item \texttt{pattern = re.compile(r`[0-9]+)}
\item \texttt{pattern.sub(``-'',``order0, order1 order13'')}
\item \texttt{`order- order- order-'}
\end{itemize}
\end{ej}
\end{frame}

\begin{frame}{Regular expressions in Python}
What is a MatchObject? \pause
\begin{block}{}
\begin{itemize}
\item An object that represents the matched pattern.
\item We will get one every time you execute one of these operations:
\begin{itemize}
\item match
\item search
\item finditer
\end{itemize}
\item Provides a set of operations for working with the captured groups.
\end{itemize}
\end{block}
\end{frame}

\begin{frame}{Regular expressions in Python}
Which are the main operations for a MatchObject? \pause
\begin{block}{}
\begin{itemize}
\item group
\item groups
\item groupdict
\item start
\item end
\item span
\end{itemize}
\end{block}
\end{frame}

\begin{frame}{Regular expressions in Python}
What is the group operation? \pause
\begin{block}{}
\begin{itemize}
\item Gives the subgroups of the match. 
\item If invoked with no arguments or zero, returns the entire match. 
\item If one or more group identifiers are passed, the corresponding groups' \, matches will be returned.
\end{itemize}
\end{block}
\end{frame}

\begin{frame}{Regular expressions in Python}
\begin{ej}
\begin{itemize}
\item \texttt{pattern = re.compile(``(\textbackslash w+) (\textbackslash w+)")}
\item \texttt{match = pattern.search(``Hello world'')}
\item \texttt{match.group()} $\rightarrow$ \texttt{`Hello world'}
\item \texttt{match.group(0)} $\rightarrow$ \texttt{`Hello world'}
\item \texttt{match.group(1)} $\rightarrow$ \texttt{`Hello'}
\item \texttt{match.group(2)} $\rightarrow$ \texttt{`world'}
\item \texttt{match.group(0,2)} $\rightarrow$ \texttt{(`Hello world',`world')}
\end{itemize}
\end{ej}
\end{frame}

\begin{frame}{Regular expressions in Python}
\begin{ej}
\begin{itemize}
\item Groups can be named.
\item If the pattern has named groups, they can be accessed using the names or the index:
\item \texttt{pattern = re.compile(r``(?P<first>\textbackslash w+) (?P<second>\textbackslash w+)")}
\item \texttt{match = pattern.search(``Hello world'')}
\item \texttt{match.group(`first')} $\rightarrow$ \texttt{`Hello'}
\item \texttt{match.group(1)} $\rightarrow$ \texttt{`Hello'}
\item \texttt{match.group(0,`first',2)} $\rightarrow$ \texttt{(`Hello world',`Hello',`world')}
\end{itemize}
\end{ej}
\end{frame}

\begin{frame}{Regular expressions in Python}
What is the groups operation? \pause
\begin{block}{}
\begin{itemize}
\item It returns a tuple with all the subgroups in the match instead of giving one or some of the groups. 
\end{itemize}
\end{block}
\end{frame}

\begin{frame}{Regular expressions in Python}
\begin{ej}
\begin{itemize}
\item \texttt{pattern = re.compile(``(\textbackslash w+) (\textbackslash w+)")}
\item \texttt{match = pattern.search(``Hello world'')}
\item \texttt{match.groups()} $\rightarrow$ \texttt{(`Hello',`world')}
\end{itemize}
\end{ej}
\end{frame}

\begin{frame}{Regular expressions in Python}
What is the groupdict operation? \pause
\begin{block}{}
\begin{itemize}
\item It is used in the cases where named groups have been used. 
\item It will return a dictionary with all the groups that were found. 
\item \texttt{pattern = re.compile(r``(?P<first>\textbackslash w+) (?P<second>\textbackslash w+)")}
\item \texttt{match = pattern.search(``Hello world'')}
\item \texttt{\{`first': `Hello', `second': `world'\}}
\item If there aren't named groups, then it returns an empty dictionary.
\end{itemize}
\end{block}
\end{frame}

\begin{frame}{Regular expressions in Python}
What about start, end and span operations? \pause
\begin{block}{}
\begin{itemize}
\item start: returns the index where the pattern matched. 
\item end: returns the end of the substring matched by the group.
\item span: gives a tuple with the values from start and end.
\end{itemize}
\end{block}
\end{frame}

\subsection{Exercises in Python}

\end{document}
